\section{Building and verifying the formalization}\label{app:verify}

The Lean project is the directory \leanfile{lean/} next to this paper. It uses Lean and Mathlib, both at version v4.34.1
(\leanfile{lean-toolchain}, \leanfile{lakefile.toml}, \leanfile{lake-manifest.json}). It requires, through relative paths,
the Lean library of the fourth paper~\cite{FourthPaper2026}, in the directory
\leanfile{Curved-Motion-Depth-Reduction-and-Storage-Codes-for-Matrix-Transposition-on-Tapes/lean}, and that of its circular
memo~\cite{CircularMemo2026}, in the subdirectory \leanfile{memo/lean} of the same paper. The fourth paper's library in turn
imports the second paper's library, which imports that of the first~\cite{FirstPaper2026,SecondPaper2026}, and the
repository's shared library \texttt{lean-machine}. Keep the repository's directory layout, so that these sibling directories
stay in place.

\paragraph{Build.} With \texttt{elan} installed, run in \leanfile{lean/}:
\begin{center}
\texttt{lake exe cache get \&\& lake build}
\end{center}
The first command downloads the compiled Mathlib; the second compiles the library and the parts of the earlier libraries
that it uses. Before the first build, fetch the second paper's pinned external dependency as its README describes.

\paragraph{Axioms.} The file \leanfile{PaperAxioms.lean}, generated from the \LaTeX{} sources of this paper by
\leanfile{tools/lean_index.py}, runs \texttt{\#print axioms} on every declaration of this library named here
(\cref{app:index}):
\begin{center}
\texttt{lake env lean PaperAxioms.lean}
\end{center}
Every declaration depends at most on \texttt{propext}, \texttt{Classical.choice} and \texttt{Quot.sound}. No Lean file of this
library uses \texttt{sorry}, \texttt{admit} or \texttt{native\_decide}, or declares an \texttt{axiom}. Running
\texttt{python3 tools/lean\_index.py --index-only} in the paper's directory regenerates only the index of
\cref{app:index}.

\paragraph{Where things are.} \leanfile{NestedElimination.lean} (Lemma~E), \leanfile{Decomposition.lean} (contraction,
piece certificates and Theorem~D), \leanfile{Positional.lean} (co-supply), \leanfile{Pieces.lean} (fewer arcs, greedy
certificates, Lemma~$\Bdir$), \leanfile{Corollaries.lean} (global certificates, the pull-back, $2$-dicolourable relations and
circular chromatic number below three), \leanfile{N3.lean} (Theorem~N3 and the shortest-cycle decomposition of blocks),
\leanfile{U3.lean} (Theorem~U3), \leanfile{TheoremT.lean} (Theorem~T and the block form of Corollary~$\Tint$),
\leanfile{Blocks.lean} (internal certificates and chordless subcubic blocks), \leanfile{PathCert.lean} (Lemma~S and
Theorem~PC), \leanfile{Circulant.lean} (positional certificates and Proposition~CC), \leanfile{Cayley.lean}
(Corollary~CC), \leanfile{Barrier.lean} (the barrier bound), \leanfile{Matching.lean} (fractional matchings, the fractional
Tutte--Berge formula and the guessing number), \leanfile{IndexCoding.lean} (index codes and the broadcast rate), and
\leanfile{Examples.lean} (small sanity checks).

\paragraph{Computations.} The numerical remarks of \cref{sec:limits} (\cref{rem:X15,rem:sweeps,rem:petersen}) and the
numerical check of the folded-cube formula in \cref{rem:foldcube} are reproduced by \leanfile{tools/verify_remarks.py}
(plain Python~3, run from the paper's directory as \texttt{python3 tools/verify\_remarks.py}), described in
\leanfile{tools/README.md}; the script also checks \cref{prop:no-private} by exhaustion, and \cref{ex:D9} and the global
certificate before it step by step. They are not used in any proof, and the Lean library does not depend on them.
